\chapter{2003年天津卷（文）}
\section{单选题}
\begin{problem}
不等式 \(\sqrt{4x - x^2} < x\) 的解集是（\(\quad\)）

\choices
    {\((0,2)\)}
    {\((2,+\infty)\)}
    {\((2,4)\)}
    {\((-\infty,0)\cup(2,+\infty)\)}
\end{problem}

\begin{answer}
C
\end{answer}

\begin{solution}
根式有意义的条件为 \(4x-x^2\geqslant0\)，即 \(x\in[0,4]\). 又因为左边非负，而左边小于 \(x\)，所以必须有 \(x>0\). 在此条件下两边可平方，得 \(4x-x^2<x^2\)，即 \(2x(x-2)>0\)，从而 \(x>2\). 与定义域交集为 \((2,4]\). 因此严格按实数定义所得解集为 \((2,4]\)；原卷选项 C 写成 \((2,4)\)，漏掉了满足不等式的端点 \(x=4\)，按原卷给出的选项答案记为 C.
\end{solution}

\begin{problem}
已知 \(x \in \left(-\frac{\pi}{2},0\right)\)，\(\cos x = \frac{4}{5}\)，则 \(\tan 2x =\)（\(\quad\)）

\choices
    {\(\frac{7}{24}\)}
    {\(-\frac{7}{24}\)}
    {\(\frac{24}{7}\)}
    {\(-\frac{24}{7}\)}
\end{problem}

\begin{answer}
D
\end{answer}

\begin{solution}
由 \(x\in\left(-\frac{\pi}{2},0\right)\) 且 \(\cos x=\frac45\)，可知 \(\sin x=-\frac35\)，所以 \(\tan x=-\frac34\). 于是
\[
\tan 2x=\frac{2\tan x}{1-\tan^2x}=\frac{2\left(-\frac34\right)}{1-\frac{9}{16}}=-\frac{24}{7}
\]
故选 D.
\end{solution}

\begin{problem}
\(\frac{1 - \sqrt{3}\i}{(\sqrt{3} + \i)^2} =\)（\(\quad\)）

\choices
    {\(\frac{1}{4} + \frac{\sqrt{3}}{4}\i\)}
    {\(-\frac{1}{4} - \frac{\sqrt{3}}{4}\i\)}
    {\(\frac{1}{2} + \frac{\sqrt{3}}{2}\i\)}
    {\(-\frac{1}{2} - \frac{\sqrt{3}}{2}\i\)}
\end{problem}

\begin{answer}
B
\end{answer}

\begin{solution}
先计算分母：
\[
(\sqrt{3}+\i)^2=2+2\sqrt{3}\i=2(1+\sqrt{3}\i)
\]
因此
\[
\frac{1-\sqrt{3}\i}{(\sqrt{3}+\i)^2}
=\frac{1-\sqrt{3}\i}{2(1+\sqrt{3}\i)}
=\frac{(1-\sqrt{3}\i)^2}{2(1+\sqrt{3}\i)(1-\sqrt{3}\i)}
=\frac{-2-2\sqrt{3}\i}{8}
=-\frac14-\frac{\sqrt{3}}4\i
\]
故选 B.
\end{solution}

\begin{problem}
已知 \(x \in \left(-\frac{\pi}{2},0\right)\)，\(\cos x = \frac{4}{5}\)，则 \(\tan 2x =\)（\(\quad\)）

\choices
    {\(\frac{7}{24}\)}
    {\(-\frac{7}{24}\)}
    {\(\frac{24}{7}\)}
    {\(-\frac{24}{7}\)}
\end{problem}

\begin{answer}
D
\end{answer}

\begin{solution}
由 \(x\in\left(-\frac{\pi}{2},0\right)\) 且 \(\cos x=\frac45\)，可知 \(\sin x=-\frac35\)，所以 \(\tan x=-\frac34\). 于是
\[
\tan 2x=\frac{2\tan x}{1-\tan^2x}=\frac{2\left(-\frac34\right)}{1-\frac{9}{16}}=-\frac{24}{7}
\]
故选 D.
\end{solution}

\begin{problem}
等差数列 \(\{a_{n}\}\) 中，已知 \(a_1 = \frac{1}{3}\)，\(a_2 + a_5 = 4\)，\(a_n = 33\)，则 \(n\) 为（\(\quad\)）

\choices
    {\(48\)}
    {\(49\)}
    {\(50\)}
    {\(51\)}
\end{problem}

\begin{answer}
C
\end{answer}

\begin{solution}
设等差数列的公差为 \(d\)，则 \(a_2=\frac13+d\)，\(a_5=\frac13+4d\). 由条件
\[
\frac13+d+\frac13+4d=4
\]
得 \(d=\frac23\). 所以
\[
a_n=\frac13+(n-1)\frac23=\frac{2n-1}{3}
\]
由 \(a_n=33\) 得 \(\frac{2n-1}{3}=33\)，即 \(n=50\). 故选 C.
\end{solution}

\begin{problem}
双曲线虚轴的一个端点为 \(M\)，两个焦点为 \(F_{1}\)，\(F_{2}\)，\(\angle F_{1}MF_{2} = 120^{\circ}\)，则双曲线的离心率为（\(\quad\)）

\choices
    {\(\sqrt{3}\)}
    {\(\frac{\sqrt{6}}{2}\)}
    {\(\frac{\sqrt{6}}{3}\)}
    {\(\frac{\sqrt{3}}{3}\)}
\end{problem}

\begin{answer}
B
\end{answer}

\begin{solution}
设双曲线的实半轴、虚半轴和半焦距分别为 \(a\)、\(b\)、\(c\)，则 \(c^2=a^2+b^2\). 取虚轴端点 \(M\) 及焦点 \(F_1\)，\(F_2\)，有
\[
MF_1=MF_2=\sqrt{b^2+c^2}
\]
在三角形 \(F_1MF_2\) 中，\(F_1F_2=2c\)，且 \(\angle F_1MF_2=120^\circ\)，由余弦定理得
\[
4c^2=2(b^2+c^2)-2(b^2+c^2)\cos120^\circ=3(b^2+c^2)
\]
代入 \(c^2=a^2+b^2\)，得到 \(4(a^2+b^2)=3(a^2+2b^2)\)，即 \(a^2=2b^2\). 于是
\[
e=\frac ca=\frac{\sqrt{a^2+b^2}}{a}=\sqrt{\frac32}=\frac{\sqrt6}{2}
\]
故选 B.
\end{solution}

\begin{problem}
设函数 \(f(x) = \begin{cases}
2^{-x} - 1, & x \leqslant 0,\\
x^{\frac{1}{2}}, & x > 0,
\end{cases}\) 若 \(f(x_0) > 1\)，则 \(x_0\) 的取值范围是（\(\quad\)）

\choices
    {\((-1,1)\)}
    {\((-1,+\infty)\)}
    {\((-\infty,-2)\cup(0,+\infty)\)}
    {\((-\infty,-1)\cup(1,+\infty)\)}
\end{problem}

\begin{answer}
D
\end{answer}

\begin{solution}
当 \(x_0\leqslant0\) 时，\(f(x_0)>1\) 等价于 \(2^{-x_0}-1>1\)，即 \(2^{-x_0}>2\)，从而 \(-x_0>1\)，所以 \(x_0<-1\). 当 \(x_0>0\) 时，\(f(x_0)>1\) 等价于 \(\sqrt{x_0}>1\)，从而 \(x_0>1\). 合并两种情形，得
\[
x_0\in(-\infty,-1)\cup(1,+\infty)
\]
故选 D.
\end{solution}

\begin{problem}
\(O\) 是平面上一定点，\(A\)，\(B\)，\(C\) 是平面上不共线的三个点，动点 \(P\) 满足 \(\overrightarrow{OP} = \overrightarrow{OA} + \lambda \left(\frac{\overrightarrow{AB}}{|\overrightarrow{AB}|} + \frac{\overrightarrow{AC}}{|\overrightarrow{AC}|}\right)\)，\(\lambda \in [0,+\infty)\)，则 \(P\) 的轨迹一定通过 \(\triangle ABC\) 的（\(\quad\)）

\choices
    {外心}
    {内心}
    {重心}
    {垂心}
\end{problem}

\begin{answer}
B
\end{answer}

\begin{solution}
记 \(\bs{u}=\frac{\overrightarrow{AB}}{|\overrightarrow{AB}|}\)、\(\bs{v}=\frac{\overrightarrow{AC}}{|\overrightarrow{AC}|}\)，则 \(|\bs{u}|=|\bs{v}|=1\). 由于 \(A\)，\(B\)，\(C\) 不共线，\(\bs{u}+\bs{v}\ne\bs{0}\)，且
\[
(\bs{u}+\bs{v})\cdot\bs{u}=1+\bs{u}\cdot\bs{v}=(\bs{u}+\bs{v})\cdot\bs{v}
\]
所以向量 \(\bs{u}+\bs{v}\) 与 \(\bs{u}\)、\(\bs{v}\) 的夹角相等，正好沿顶点 \(A\) 的内角平分线方向. 由题设
\(\overrightarrow{AP}=\overrightarrow{OP}-\overrightarrow{OA}=\lambda(\bs{u}+\bs{v})\)，\(\lambda\geqslant0\)
故 \(P\) 的轨迹是从 \(A\) 出发沿该内角平分线方向的射线. 三角形的内心在三个内角平分线的交点上，且位于从 \(A\) 出发的内角平分线射线上，所以该轨迹必经过内心，故选 B.
\end{solution}

\begin{problem}
函数 \(y = \ln \frac{x + 1}{x - 1}\)，\(x \in (1,+\infty)\) 的反函数为（\(\quad\)）

\choices
    {\(y = \frac{\e^x - 1}{\e^x + 1}, x \in (0,+\infty)\)}
    {\(y = \frac{\e^x + 1}{\e^x - 1}\)，\(x \in (0,+\infty)\)}
    {\(y = \frac{\e^x - 1}{\e^x + 1}\)，\(x \in (-\infty,0)\)}
    {\(y = \frac{\e^x + 1}{\e^x - 1}\)，\(x \in (-\infty,0)\)}
\end{problem}

\begin{answer}
B
\end{answer}

\begin{solution}
设原函数的函数值为 \(y\)，则
\[
\e^y=\frac{x+1}{x-1}
\]
整理得 \(\e^y x-\e^y=x+1\)，所以
\[
x=\frac{\e^y+1}{\e^y-1}
\]
当 \(x>1\) 时，\(\frac{x+1}{x-1}>1\)，因此 \(y>0\). 交换自变量和函数值，得到反函数
\(y=\frac{\e^x+1}{\e^x-1}\)，\(x\in(0,+\infty)\).
故选 B.
\end{solution}

\begin{problem}
棱长为 \(a\) 的正方体中，连结相邻面的中心，以这些线段为棱的八面体的体积为（\(\quad\)）

\choices
    {\(\frac{a^3}{3}\)}
    {\(\frac{a^3}{4}\)}
    {\(\frac{a^3}{6}\)}
    {\(\frac{a^3}{12}\)}
\end{problem}

\begin{answer}
C
\end{answer}

\begin{solution}
以正方体中心为原点建立直角坐标系，则六个面的中心为
\(\left(\pm\frac a2,0,0\right)\)，\(\left(0,\pm\frac a2,0\right)\)，\(\left(0,0,\pm\frac a2\right)\).
所得八面体可分成上、下两个全等的四棱锥. 以 \(z=0\) 平面内的正方形为共同底面，其两条对角线均为 \(a\)，所以底面积为 \(\frac{a^2}{2}\)，每个四棱锥的高为 \(\frac a2\). 因此八面体体积为
\[
2\cdot\frac13\cdot\frac{a^2}{2}\cdot\frac a2=\frac{a^3}{6}
\]
故选 C.
\end{solution}

\begin{problem}
设 \(a > 0\)，\(f(x) = ax^2 + bx + c\)，曲线 \(y = f(x)\) 在点 \(P(x_0, f(x_0))\) 处切线的倾斜角的取值范围为 \(\left[0, \frac{\pi}{4}\right]\)，则 \(P\) 到曲线 \(y = f(x)\) 对称轴距离的取值范围为（\(\quad\)）

\choices
    {\(\left[0, \frac{1}{a}\right]\)}
    {\(\left[0, \frac{1}{2a}\right]\)}
    {\(\left[0, \left|\frac{b}{2a}\right|\right]\)}
    {\(\left[0, \left|\frac{b-1}{2a}\right|\right]\)}
\end{problem}

\begin{answer}
B
\end{answer}

\begin{solution}
曲线在 \(P(x_0,f(x_0))\) 处的切线斜率为
\[
m=f'(x_0)=2ax_0+b
\]
由切线倾斜角 \(\alpha\in\left[0,\frac{\pi}{4}\right]\)，得 \(m=\tan\alpha\in[0,1]\). 抛物线的对称轴为 \(x=-\frac{b}{2a}\)，所以 \(P\) 到对称轴的距离为
\[
\left|x_0+\frac{b}{2a}\right|
=\left|\frac{2ax_0+b}{2a}\right|
=\frac{m}{2a}
\]
因为 \(m\in[0,1]\) 且 \(a>0\)，该距离的取值范围为 \(\left[0,\frac{1}{2a}\right]\). 故选 B.
\end{solution}

\begin{problem}
已知长方形的四个顶点 \(A(0,0)\)，\(B(2,0)\)，\(C(2,1)\) 和 \(D(0,1)\)，一质点从 \(AB\) 的中点 \(P_0\) 沿与 \(AB\) 的夹角为 \(\theta\) 的方向射到 \(BC\) 上的点 \(P_1\) 后，依次反射到 \(CD\)，\(DA\) 和 \(AB\) 上的点 \(P_2\)，\(P_3\) 和 \(P_4\)（入射角等于反射角），设 \(P_4\) 的坐标为 \((x_4,0)\)，若 \(1 < x_4 < 2\)，则 \(\tan \theta\) 的取值范围是（\(\quad\)）

\choices
    {\(\left(\frac{1}{3},1\right)\)}
    {\(\left(\frac{1}{3},\frac{2}{3}\right)\)}
    {\(\left(\frac{2}{5},\frac{1}{2}\right)\)}
    {\(\left(\frac{2}{5},\frac{3}{5}\right)\)}
\end{problem}

\begin{answer}
C
\end{answer}

\begin{solution}
设 \(t=\tan\theta\). 质点从 \(P_0=(1,0)\) 射向右上方，第一次到达 \(BC\) 时
\[
P_1=(2,t)
\]
故 \(0<t<1\). 在 \(x=2\) 处反射后向左上方运动，到达 \(CD\) 时的点为
\[
P_2=\left(2-\frac{1-t}{t},1\right)=\left(3-\frac1t,1\right)
\]
由 \(0<x_{P_2}<2\)，得 \(\frac13<t<1\). 在 \(CD\) 处反射后向左下方运动，到达 \(DA\) 时
\[
P_3=\left(0,1-\left(3-\frac1t\right)t\right)=(0,2-3t)
\]
由 \(0<2-3t<1\)，得 \(\frac13<t<\frac23\). 在 \(DA\) 处反射后向右下方运动，到达 \(AB\) 时的横坐标为
\[
x_4=\frac{2-3t}{t}=\frac2t-3
\]
最后的条件 \(1<x_4<2\) 等价于
\[
1<\frac2t-3<2
\quad\Longleftrightarrow\quad
\frac25<t<\frac12
\]
该区间已包含在前面的命中条件内，所以 \(\tan\theta\in\left(\frac25,\frac12\right)\)，故选 C.
\end{solution}

\begin{problem}
一个四面体的所有棱长都为 \(\sqrt{2}\)，四个顶点在同一球面上，则此球的表面积为（\(\quad\)）

\choices
    {\(3\pi\)}
    {\(4\pi\)}
    {\(3\sqrt{3}\pi\)}
    {\(6\pi\)}
\end{problem}

\begin{answer}
A
\end{answer}

\begin{solution}
这是棱长 \(s=\sqrt2\) 的正四面体. 设底面为正三角形，底面中心为 \(G\)，顶点为 \(A\)，外接球球心为 \(O\)，外接球半径为 \(R\). 底面中心到任一底面顶点的距离为 \(\frac{s}{\sqrt3}\)，所以正四面体的高为
\[
h=AG=\sqrt{s^2-\frac{s^2}{3}}=s\sqrt{\frac23}
\]
球心在高 \(AG\) 上，并且 \(OG=h-R\). 由 \(OA=OB=R\) 以及 \(GB=\frac{s}{\sqrt3}\)，有
\[
R^2=(h-R)^2+\frac{s^2}{3}
\]
化简得 \(2hR=h^2+\frac{s^2}{3}=s^2\)，从而
\[
R=\frac{s^2}{2h}=\frac{s\sqrt6}{4}=\frac{\sqrt3}{2}
\]
故球的表面积为
\[
4\pi R^2=4\pi\cdot\frac34=3\pi
\]
故选 A.
\end{solution}

\section{填空题}
\begin{problem}
\(\left(x^{2} - \frac{1}{2x}\right)^{9}\) 的展开式中 \(x^{9}\) 系数是\fillinblank{}.
\end{problem}

\begin{answer}
\(-\frac{21}{2}\)
\end{answer}

\begin{solution}
展开式的通项为
\(T_{k+1}=\mathrm C_9^k(x^2)^{9-k}\left(-\frac{1}{2x}\right)^k =\mathrm C_9^k\frac{(-1)^k}{2^k}x^{18-3k}\)，\(k=0\)，\(1\)，\(\ldots\)，\(9\).
要得到 \(x^9\) 项，必须满足 \(18-3k=9\)，所以 \(k=3\). 因此该项的系数为
\[
\mathrm C_9^3\frac{(-1)^3}{2^3}
=-\frac{84}{8}
=-\frac{21}{2}
\]
\end{solution}

\begin{problem}
某公司生产三种型号的轿车，产量分别为 \(1200\) 辆，\(6000\) 辆和 \(2000\) 辆. 为检验该公司的产品质量，现用分层抽样的方法抽取 \(46\) 辆进行检验，这三种型号的轿车依次应抽取\fillinblank{}，\fillinblank{}，\fillinblank{} 辆.
\end{problem}

\begin{answer}
\(6\)，\(30\)，\(10\)
\end{answer}

\begin{solution}
三种型号的轿车总数为 \(1200+6000+2000=9200\)，产量之比为
\[
1200:6000:2000=3:15:5
\]
总样本量 \(46=2\times(3+15+5)\)，所以按比例抽取的三层样本数分别为
\(2\times3=6\)，\(2\times15=30\)，\(2\times5=10\).
\end{solution}

\begin{problem}
在平面几何里，有勾股定理：“设 \(\triangle ABC\) 的两边 \(AB\)，\(AC\) 互相垂直，则 \(AB^2 + AC^2 = BC^2\). ”拓展到空间，类比平面几何的勾股定理，研究三棱锥的侧面面积与底面面积间的关系，可以得出的正确结论是：“设三棱锥 \(A - BCD\) 的三个侧面 \(ABC\)，\(ACD\)，\(ADB\) 两两互相垂直，则\fillinblank{}.”
\end{problem}

\begin{answer}
\(S_{ABC}^{2} + S_{ACD}^{2} + S_{ADB}^{2} = S_{BCD}^{2}\)
\end{answer}

\begin{solution}
记三个侧面 \(ABC\)、\(ACD\)、\(ADB\) 所在平面的单位法向量分别为 \(\bs{n}_1\)、\(\bs{n}_2\)、\(\bs{n}_3\). 由三个平面两两垂直，\(\bs{n}_1\)、\(\bs{n}_2\)、\(\bs{n}_3\) 两两垂直. 又 \(AC\)、\(AB\)、\(AD\) 分别是这三个平面的两两交线，故它们分别平行于 \(\bs{n}_1\times\bs{n}_2\)、\(\bs{n}_1\times\bs{n}_3\)、\(\bs{n}_2\times\bs{n}_3\)，从而 \(AB\)、\(AC\)、\(AD\) 两两垂直. 设
\(AB=p\)，\(AC=q\)，\(AD=r\).
则三个侧面面积分别为
\(S_{ABC}=\frac{pq}{2}\)，\(S_{ACD}=\frac{qr}{2}\)，\(S_{ADB}=\frac{rp}{2}\).
建立以 \(AB\)、\(AC\)、\(AD\) 为坐标轴的直角坐标系，取 \(A\) 为原点，则 \(B=(p,0,0)\)、\(C=(0,q,0)\)、\(D=(0,0,r)\). 于是
\(\overrightarrow{BC}=(-p,q,0)\)，\(\overrightarrow{BD}=(-p,0,r)\)
从而
\[
S_{BCD}
=\frac12\left|\overrightarrow{BC}\times\overrightarrow{BD}\right|
=\frac12\sqrt{p^2q^2+q^2r^2+r^2p^2}
\]
两边平方即得
\[
S_{BCD}^{2}
=\frac{p^2q^2+q^2r^2+r^2p^2}{4}
=S_{ABC}^{2}+S_{ACD}^{2}+S_{ADB}^{2}
\]
\end{solution}

\begin{problem}
将 \(3\) 种作物种植在如图 \(5\) 块试验田里，每块种植一种作物且相邻的试验田不能种植同一种作物，不同的种植方法共有\fillinblank{} 种.（以数字作答）

\begin{center}
\begin{tabular}{|c|c|c|c|c|}
\hline
\rule{2em}{0pt}\rule{0pt}{2em} &
\rule{2em}{0pt} &
\rule{2em}{0pt} &
\rule{2em}{0pt} &
\rule{2em}{0pt} \\
\hline
\end{tabular}
\end{center}
\end{problem}

\begin{answer}
42
\end{answer}

\begin{solution}
因为题目要求将 \(3\) 种作物都种植出来，先不考虑三种作物都出现的条件. 第一块试验田有 \(3\) 种选择，之后每一块都必须与前一块不同，均有 \(2\) 种选择，所以相邻不同的种植序列共有
\[
3\cdot2^4=48
\]
种.

其中没有使用三种作物的序列只能使用其中任意两种作物. 选取这两种作物有 \(\mathrm C_3^2=3\) 种，选定后第一块有 \(2\) 种选择，后面各块的作物被相邻不同条件唯一确定，故这类序列有
\[
\mathrm C_3^2\cdot2=6
\]
种. 由于只使用一种作物不可能满足相邻不同条件，所求种植方法数为
\[
48-6=42
\]
\end{solution}

\section{解答题}
\begin{problem}
已知正四棱柱 \(ABCD - A_{1}B_{1}C_{1}D_{1}\)，\(AB = 1\)，\(AA_{1} = 2\)，\(E\) 为 \(CC_{1}\) 中点，\(F\) 为 \(BD_{1}\) 中点.

\begin{enumerate}
\item 证明： \(EF\) 为 \(BD_{1}\) 与 \(CC_{1}\) 的公垂线；
\item 求点 \(D_{1}\) 到面 \(BDE\) 的距离.
\end{enumerate}
\bitmapfigure[max width=4.8cm]{img/2003/tianjin_liberal/q18_fig1.png}

\begin{solution}
\begin{enumerate}
\item 取 \(A\) 为原点，\(\overrightarrow{AB}\)、\(\overrightarrow{AD}\)、\(\overrightarrow{AA_{1}}\) 分别为 \(x\)、\(y\)、\(z\) 轴的正向，建立直角坐标系，则
\(B(1,0,0)\)、\(C(1,1,0)\)、\(D(0,1,0)\)、\(D_{1}(0,1,2)\)、\(C_{1}(1,1,2)\)，从而
\(E(1,1,1)\)、\(F\left(\frac{1}{2},\frac{1}{2},1\right)\).

于是 \(\overrightarrow{EF}=\left(-\frac{1}{2},-\frac{1}{2},0\right)\)、\(\overrightarrow{BD_{1}}=(-1,1,2)\)、\(\overrightarrow{CC_{1}}=(0,0,2)\)，故
\(\overrightarrow{EF}\cdot\overrightarrow{BD_{1}}=0\)，\(\overrightarrow{EF}\cdot\overrightarrow{CC_{1}}=0\). 又 \(F\in BD_{1}\)，\(E\in CC_{1}\)，所以 \(EF\) 同时垂直于 \(BD_{1}\) 和 \(CC_{1}\)，即 \(EF\) 为它们的公垂线.

\item \(\overrightarrow{BD}=(-1,1,0)\)，\(\overrightarrow{BE}=(0,1,1)\). 取平面 \(BDE\) 的一个法向量为
\[
\bs{n}=\overrightarrow{BD}\times\overrightarrow{BE}=(1,1,-1)
\]
因此平面 \(BDE\) 的方程为 \(x+y-z-1=0\). 点 \(D_{1}(0,1,2)\) 到该平面的距离为
\[
d=\frac{|0+1-2-1|}{\sqrt{1^{2}+1^{2}+(-1)^{2}}}=\frac{2}{\sqrt{3}}=\frac{2\sqrt{3}}{3}
\]
\end{enumerate}
\end{solution}
\end{problem}

\begin{problem}
已知抛物线 \(C_{1}: y = x^{2} + 2x\) 和 \(C_{2}: y = -x^{2} + a\)，如果直线 \(l\) 同时是 \(C_{1}\) 和 \(C_{2}\) 的切线，称 \(l\) 是 \(C_{1}\) 和 \(C_{2}\) 的公切线. 公切线上两个切点之间的线段，称为公切线段.

\begin{enumerate}
\item \(a\) 取什么值时，\(C_1\) 和 \(C_2\) 有且仅有一条公切线？写出此公切线的方程；
\item 若 \(C_1\) 和 \(C_2\) 有两条公切线，证明相应的两条公切线段互相平分.
\end{enumerate}

\begin{solution}
设公切线 \(l\) 的方程为 \(y=kx+b\). 由于两条抛物线的切线均不可能为竖直直线，这样设不失一般性.

直线 \(l\) 与 \(C_{1}\) 联立得 \(x^{2}+(2-k)x-b=0\). 该方程有重根，所以
\[
(2-k)^{2}+4b=0
\]
即 \(b=-\frac{(k-2)^{2}}{4}\).

直线 \(l\) 与 \(C_{2}\) 联立并整理得 \(x^{2}+kx+b-a=0\). 该方程有重根，所以
\[
k^{2}-4(b-a)=0
\]
即 \(b=a+\frac{k^{2}}{4}\).

消去 \(b\)，得到公切线的斜率 \(k\) 满足
\[
k^{2}-2k+2+2a=0
\]

\begin{enumerate}
\item 公切线的条数由上述关于 \(k\) 的二次方程的实根个数决定. 有且仅有一条公切线时，判别式为零，即
\[
(-2)^{2}-4(2+2a)=0
\]
所以 \(a=-\frac{1}{2}\)，且 \(k=1\). 此时 \(b=-\frac{(1-2)^{2}}{4}=-\frac{1}{4}\)，故公切线方程为
\[
y=x-\frac{1}{4}
\]

\item 设两条公切线的斜率分别为 \(k_{1}\)、\(k_{2}\). 对任意一个公切线斜率 \(k\)，它与 \(C_{1}\) 的切点横坐标为二次方程 \(x^{2}+(2-k)x-b=0\) 的重根，即 \(x=\frac{k-2}{2}\)，所以切点为
\[
T(k)=\left(\frac{k-2}{2},\frac{k^{2}-4}{4}\right)
\]
它与 \(C_{2}\) 的切点横坐标为二次方程 \(x^{2}+kx+b-a=0\) 的重根，即 \(x=-\frac{k}{2}\)，所以切点为
\[
S(k)=\left(-\frac{k}{2},a-\frac{k^{2}}{4}\right)
\]

公切线段 \(T(k)S(k)\) 的中点为
\[
\left(\frac{\frac{k-2}{2}-\frac{k}{2}}{2},\frac{\frac{k^{2}-4}{4}+a-\frac{k^{2}}{4}}{2}\right)=\left(-\frac{1}{2},\frac{a-1}{2}\right)
\]
与 \(k\) 无关. 因而对应于 \(k_{1}\)、\(k_{2}\) 的两条公切线段有同一个中点，即两条公切线段互相平分.
\end{enumerate}
\end{solution}
\end{problem}

\begin{problem}
已知数列 \(\{a_{n}\}\) 满足 \(a_{1} = 1\)，\(a_{n} = 3^{n - 1} + a_{n - 1}\)（\(n \geqslant 2\)）.

\begin{enumerate}
\item 求 \(a_{2}\)，\(a_{3}\).
\item 证明： \(a_{n} = \frac{3^{n} - 1}{2}\).
\end{enumerate}

\begin{solution}
\begin{enumerate}
\item 由递推关系，\(a_{2}=3^{1}+a_{1}=3+1=4\)，\(a_{3}=3^{2}+a_{2}=9+4=13\).

\item 当 \(n=1\) 时，\(a_{1}=1=\frac{3-1}{2}\)，结论成立.

假设当 \(n=k\) 时结论成立，即 \(a_{k}=\frac{3^{k}-1}{2}\). 则由递推关系
\[
a_{k+1}=3^{k}+a_{k}=3^{k}+\frac{3^{k}-1}{2}=\frac{3^{k+1}-1}{2}
\]
因此，当 \(n=k+1\) 时结论也成立. 由数学归纳法，\(a_{n}=\frac{3^{n}-1}{2}\) 对一切正整数 \(n\) 成立.
\end{enumerate}
\end{solution}
\end{problem}

\begin{problem}
有三种产品，合格率分别为 \(0.90\)，\(0.95\) 和 \(0.95\)，各抽取一件进行检验.

\begin{enumerate}
\item 求恰有一件不合格的概率；
\item 求至少有两件不合格的概率.（精确到 \(0.001\)）
\end{enumerate}

\begin{solution}
设三件产品不合格的概率分别为 \(0.10\)、\(0.05\)、\(0.05\)，各件产品的检验结果相互独立.

\begin{enumerate}
\item 恰有一件不合格的概率为
\[
0.10\times0.95\times0.95+0.90\times0.05\times0.95+0.90\times0.95\times0.05=0.17575\approx0.176
\]

\item 恰有两件不合格的概率为
\[
0.10\times0.05\times0.95+0.10\times0.95\times0.05+0.90\times0.05\times0.05=0.01175
\]
三件均不合格的概率为
\[
0.10\times0.05\times0.05=0.00025
\]
所以至少有两件不合格的概率为 \(0.01175+0.00025=0.012\).
\end{enumerate}
\end{solution}
\end{problem}

\begin{problem}
已知函数 \(f(x) = \sin (\omega x + \varphi)\)（\(\omega > 0\)，\(0 \leqslant \varphi \leqslant \pi\)）是 \(\R\) 上的偶函数，其图象关于点 \(M\left(\frac{3\pi}{4}, 0\right)\) 对称，且在区间 \(\left[0, \frac{\pi}{2}\right]\) 上是单调函数. 求 \(\omega\) 和 \(\varphi\) 的值.

\begin{solution}
由于 \(f(x)\) 是偶函数，对任意 \(x\) 有
\[
0=f(x)-f(-x)=\sin(\omega x+\varphi)-\sin(-\omega x+\varphi)=2\cos\varphi\sin\omega x
\]
又 \(\omega>0\)，故 \(\cos\varphi=0\). 结合 \(0\leqslant\varphi\leqslant\pi\)，得 \(\varphi=\frac{\pi}{2}\)，从而 \(f(x)=\cos\omega x\).

图象关于 \(M\left(\frac{3\pi}{4},0\right)\) 对称，故对任意 \(x\) 有
\[
\cos\left(\omega\left(\frac{3\pi}{2}-x\right)\right)=-\cos\omega x
\]
令 \(t=\omega x\)，则上式对任意 \(t\in\R\) 成立. 由
\[
\cos\left(\frac{3\pi\omega}{2}-t\right)=\cos\frac{3\pi\omega}{2}\cos t+\sin\frac{3\pi\omega}{2}\sin t
\]
可得
\(\cos\frac{3\pi\omega}{2}=-1\)，\(\sin\frac{3\pi\omega}{2}=0\)
所以 \(\frac{3\pi\omega}{2}=(2m+1)\pi\)，即 \(\omega=\frac{2(2m+1)}{3}\)，其中 \(m\) 为非负整数.

在区间 \(\left[0,\frac{\pi}{2}\right]\) 上，\(\omega x\) 的取值范围为 \(\left[0,\frac{(2m+1)\pi}{3}\right]\). 当 \(m=0\) 或 \(m=1\) 时，\(\cos\omega x\) 在该区间上单调递减；当 \(m\geqslant2\) 时，\(\omega x\) 在该区间内越过 \(\pi\)，函数先减后增，不满足单调性. 因此
\[
(\omega,\varphi)=\left(\frac{2}{3},\frac{\pi}{2}\right)\quad\text{或}\quad\left(2,\frac{\pi}{2}\right)
\]
\end{solution}
\end{problem}

\begin{problem}
已知常数 \(\alpha > 0\)，向量 \(\bs{c} = (0, \alpha)\)，\(\bs{i} = (1, 0)\). 经过原点 \(O\) 以 \(\bs{c} + \lambda \bs{i}\) 为方向向量的直线与经过定点 \(A(0, \alpha)\) 以 \(\bs{i} - 2\lambda \bs{c}\) 为方向向量的直线相交于 \(P\)，其中 \(\lambda \in \R\). 试问：是否存在两个定点 \(E\)，\(F\)，使得 \(|PE| + |PF|\) 为定值. 若存在，求出 \(E\)，\(F\) 的坐标；若不存在，说明理由.

\begin{solution}
设交点 \(P\) 在第一条直线上的参数为 \(t\)，在第二条直线上的参数为 \(s\)，则
\[
P=t(\lambda,\alpha)=(s,\alpha-2\alpha\lambda s)
\]
因此 \(s=t\lambda\)，且 \(\alpha t=\alpha-2\alpha\lambda s\). 代入 \(s=t\lambda\)，得
\[
t(1+2\lambda^{2})=1
\]
从而
\[
P=\left(\frac{\lambda}{1+2\lambda^{2}},\frac{\alpha}{1+2\lambda^{2}}\right)
\]

记 \(P=(x,y)\). 由上式中的 \(y>0\) 可得 \(\lambda=\frac{\alpha x}{y}\)，再由 \(y=\frac{\alpha}{1+2\lambda^{2}}\) 消去 \(\lambda\)，得到
\[
2\alpha^{2}x^{2}+y^{2}-\alpha y=0
\]
即
\[
2\alpha^{2}x^{2}+\left(y-\frac{\alpha}{2}\right)^{2}=\frac{\alpha^{2}}{4}
\]
这是一条中心为 \(\left(0,\frac{\alpha}{2}\right)\) 的椭圆，其半轴长分别为 \(\frac{1}{2\sqrt{2}}\) 和 \(\frac{\alpha}{2}\).

当 \(0<\alpha<\frac{\sqrt{2}}{2}\) 时，椭圆的长轴在 \(x\) 轴方向，两个焦点为
\(E\left(-\frac{1}{2}\sqrt{\frac{1}{2}-\alpha^{2}},\frac{\alpha}{2}\right)\)，\(F\left(\frac{1}{2}\sqrt{\frac{1}{2}-\alpha^{2}},\frac{\alpha}{2}\right)\)
且 \(|PE|+|PF|=\frac{1}{\sqrt{2}}\).

当 \(\alpha>\frac{\sqrt{2}}{2}\) 时，椭圆的长轴在 \(y\) 轴方向，两个焦点为
\(E\left(0,\frac{\alpha-\sqrt{\alpha^{2}-\frac{1}{2}}}{2}\right)\)，\(F\left(0,\frac{\alpha+\sqrt{\alpha^{2}-\frac{1}{2}}}{2}\right)\)
且 \(|PE|+|PF|=\alpha\).

当 \(\alpha=\frac{\sqrt{2}}{2}\) 时，上式表示一个圆，两个焦点重合于圆心，不能得到两个不同的定点，因此按题意不存在符合条件的两个定点.
\end{solution}
\end{problem}
